\section{逆Gram族の支配性と二つの楕円moduli}
\label{sec:inverse-gram-dominance-v29}

Hermite曲線がnetのみに依存することから、その上の核束もnetだけで決まるとは限らない。
この節では、逆Gram構成が六次元上の一般の交代二形式の三つ組を、基底変換を許して全て実現することを示す。
その結果、Hermite曲線とは別に、そのJacobian上の自由に変化する点が現れる。
以下の支配性は標数零で成立し、幾何学的moduliに関する結論は代数閉体上で述べる。

$\mathcal H=\operatorname{Sym}^2(k^2)\otimes\operatorname{Sym}^2(k^3)$ を
十八次元のブロック対称行列空間とし、
\[
 G=\begin{pmatrix}A&C\\C&B\end{pmatrix},\qquad
 D_i=\operatorname{diag}([e_i]_\times,[e_i]_\times),\qquad
 \Omega_{ij}(G)=D_i^tG^{-1}D_j-D_j^tG^{-1}D_i
\]
と置く。$A,B,C$ は対称な $3\times3$ 行列である。
$\mathcal H^\times=\{\det G\ne0\}$ と書く。

\begin{theorem}[逆Gram構成の全微分階数]
\label{thm:inverse-gram-full-dominance-v29}
射
\begin{align*}
 \mathcal F:\mathcal H^\times\times\operatorname{GL}_6
 &\longrightarrow\bigl(\textstyle\bigwedge^2(k^6)^*\bigr)^{\oplus3},\\
 (G,g)&\longmapsto
 (g^t\Omega_{23}(G)g,g^t\Omega_{31}(G)g,g^t\Omega_{12}(G)g)
\end{align*}
は支配的である。特に、出力の三形式を取り替える $\operatorname{GL}_3$ の作用を
追加せずとも、その四十五次元の標的全体を支配する。
\end{theorem}

\begin{proof}
一点での微分の全射性を有理数体上で検証する。
$g=I_6$ とし、
\[
 A=\begin{pmatrix}0&1&0\\1&0&0\\0&0&1\end{pmatrix},\quad
 C=\begin{pmatrix}1&0&0\\0&0&2\\0&2&0\end{pmatrix},\quad
 B=\begin{pmatrix}0&0&1\\0&1&0\\1&0&0\end{pmatrix}
\]
を取る。このとき $\det G=-3$ である。
微分は
\[
 d(G^{-1})[H]=-G^{-1}HG^{-1},\qquad
 d(g^t\Omega g)[E]=E^t\Omega+\Omega E
\]
により計算する。
四十五行は $(\Omega_{23},\Omega_{31},\Omega_{12})$ の順に、
各交代行列の $i<j$ 成分を辞書式順序で並べる。
五十四列は、$A,C,B$ の順の対称成分 $0\le i\le j\le2$ を辞書式順序で十八本、
続いて $\operatorname{GL}_6$ の $E_{ij}$ を $0\le i,j\le5$ の辞書式順序で三十六本並べる。
この零始まりの列番号のうち
\[
 \begin{gathered}
 0,1,2,3,4,5,6,7,8,9,10,11,12,13,14,15,16,17,\\
 18,19,20,21,22,23,24,25,26,27,28,29,30,31,33,34,\\
 36,37,38,39,40,41,42,43,44,45,46
 \end{gathered}
\]
の四十五列を選ぶと、その小行列式は
\begin{equation}\label{eq:dominance-minor-v29}
 \frac{962948}{1162261467}
 =\frac{2^2\cdot7^2\cdot17^3}{3^{19}}\ne0
\end{equation}
である。

この有限計算は付属の
\nolinkurl{ruled_moduli_dominance_v29.py}で定義から再計算できる。
全ての行を九倍すると最大絶対成分三十六の整数行列となり、
整数の分数なし消去による行列式は
\[
 7231225761666696040510882238572188368556
 =2^2\cdot3^{71}\cdot7^2\cdot17^3
\]
である。従って微分の階数は四十五であり、標数零の微分判定により支配性を得る。
同じ点のPfaffian三次曲線は非零スカラーを除いて
\[
 2z_1^3+z_2^3+z_3^3-7z_1z_2z_3=0
\]
であり、$27\cdot2\cdot1\cdot1-7^3=-289\ne0$ なので滑らかでもある。
\end{proof}

\subsection{行列式線束のmoduli}

一般の交代二形式の三つ組は二つの相補的な被約共通等方三平面を持つ。
それらを順序付けると
\[
 \Omega(z)=\begin{pmatrix}0&M(z)\\-M(z)^t&0\end{pmatrix},\qquad
 0\to\mathcal O_{\mathbb P^2}(-2)^3\xrightarrow{M}
 \mathcal O_{\mathbb P^2}(-1)^3\to i_*\theta\to0
\]
が得られる。$E=\{\det M=0\}$ が滑らかなとき
$\theta\in\operatorname{Pic}^0(E)\setminus\{\mathcal O_E\}$ である。
左右の基底変換は $\theta$ を変えず、二平面の交換は転置により
$\theta\mapsto\theta^{-1}$ を与える。

\begin{corollary}[第二の楕円moduliとnet非因子化]
\label{cor:ruled-two-moduli-v29}
逆Gram族から得る
\[
 (E,\{\theta,\theta^{-1}\})
\]
は二次元の幾何学的moduliを支配する。
従って一般netを固定した五次元の二次曲線パラメータ空間上でも、
$\operatorname{Pic}^0(E)/\{\pm1\}$ における $\theta$ の像は非定数である。
幾何学的一般点で $\theta$ は無限位数であり、複素数体上の非常に一般の点でも同様である。
また、$\theta$ および $\mathbb P_E(\mathcal K)$ の幾何学的同型類はnetのみからは決まらない。
\end{corollary}

\begin{proof}
固定した滑らかな平面三次曲線 $E$ に対し、一次行列式表示の存在と一意性
\cite[Propositions~3.1 and~3.3]{BeauvilleDeterminantal}は、
$3\times3$ の行列 $M$ の左右の基底変換による同値類と
$\operatorname{Pic}^0(E)\setminus\{\mathcal O_E\}$ を対応させる。
従って共通等方三平面を順序付けた一般の交代二形式の三つ組には、曲線 $E$ と点 $\theta$ の
二つの連続moduliがある。順序を忘れる操作は $\theta$ と $\theta^{-1}$ を同一視する有限商である。
次数三の平面埋め込みの線束は、代数閉体上では平行移動で調整できるため、
さらに連続moduliを加えるものではない。
定理\ref{thm:inverse-gram-full-dominance-v29}から、逆Gram族もこの二次元空間を支配する。

Gram--Prym表示の基底変換に対する共変性により、この不変量は
$\mathcal R_3$ の幾何学的同型類に依存する。
一方、$E$ はnetだけから決まり、netのmoduliは一次元である。
もし一般netを固定したとき $\theta$ のKummer像も定数なら、上の対の全像は
高々一次元となり、支配性に反する。

各正整数 $n$ について $\theta^n=\mathcal O_E$ は二次元のmoduli内の真の閉条件である。
従って幾何学的一般点ではどの $n$ に対しても成立せず、
複素数体上では可算個の真閉部分の合併を除いて同じ結論が成立する。
核束は $\mathcal K=\mathcal O_E(-1)\otimes(\theta^{-1}\oplus\theta)$ なので、
その射影化は $\mathbb P_E(\mathcal O_E\oplus\theta^2)$ となる。
二乗写像は有限であるため、この線織面の同型類もnetのみからは決まらない。
\end{proof}

以上は幾何学的moduliに関する主張である。
粗モジュラー基底上に普遍楕円曲線がそのまま存在するとは仮定していない。
また無限位数の結論は、単一の非空Zariski開集合の全ての複素点に対する主張ではない。
有限安定化群の反転と共通等方平面の交換を区別した降下は、次の節で扱う。
